\label{ch:nsi-implementation}

This chapter maps every equation of \cref{ch:nsi-theory} onto the concrete
\texttt{tpeanuts} implementation, module by module. There is no separate
\texttt{core.BSM} Hamiltonian-assembly module: \pyfunc{hamiltonian\_reduced}
and \pyfunc{hamiltonian\_flavour} in \texttt{core.common.hamiltonian}
already handle NSI, the 3+1 sterile extension, or both simultaneously,
dispatching purely on which fields of \pyclass{OscillationParameters} are
populated.

\section{\texttt{NSIConfig}: the parameter container}
\label{sec:impl-nsiconfig}

\modulehead{core/BSM/bsm\_nsi.py}{frozen dataclass storing every NSI
coupling and building the corresponding complex tensors.}

\begin{implbox}
\pyclass{NSIConfig} stores the electron/proton-normalized block $\epsilon$
as nine real scalar fields (three real diagonal entries plus three complex
off-diagonal entries decomposed into real/imaginary parts:
\code{eps\_ee}, \code{eps\_mumu}, \code{eps\_tautau}, \code{eps\_emu\_re/\_im},
\code{eps\_etau\_re/\_im}, \code{eps\_mutau\_re/\_im}), and the
neutron-normalized block $\epsilon^n$ (\cref{eq:epsilon-composition}) as the
mirrored nine fields \code{eps\_ee\_n}, \dots, \code{eps\_mutau\_n\_im}.
Every field defaults to \code{0.0}, so $\epsilon^n=0$ (no composition
dependence) reproduces every pre-existing single-matrix scenario exactly.

\pyfunc{\_\_post\_init\_\_} recomputes \emph{both} derived complex $3\times3$
tensors, \code{epsilon} and \code{epsilon\_n}, from their respective nine
scalar fields on \emph{every} construction -- including every
\pyfunc{dataclasses.replace(...)} call, since that re-invokes
\pyfunc{\_\_init\_\_}/\pyfunc{\_\_post\_init\_\_}. This makes the eighteen
scalar fields the single source of truth: tweaking one field via
\pyfunc{replace} can never leave a stale \code{epsilon}/\code{epsilon\_n}
behind.
\end{implbox}

\begin{implbox}[title={Differentiable tensor construction}]
Both $3\times3$ Hermitian matrices are assembled by the shared helper
\pyfunc{\_hermitian\_3x3}, built entirely from \code{torch} operations
(\code{torch.complex}, \code{torch.stack}) rather than Python's
\code{complex()} builtin or \code{torch.tensor(nested\_python\_list)}:
\begin{lstlisting}[style=tpython]
def _hermitian_3x3(ee, mumu, tautau, emu_re, emu_im,
                    etau_re, etau_im, mutau_re, mutau_im,
                    *, device, real_dtype):
    def _rt(x):
        return as_tensor(x, dtype=real_dtype, device=resolved_device)

    zero = torch.zeros((), dtype=real_dtype, device=resolved_device)
    ee_c, mumu_c, tautau_c = (torch.complex(_rt(v), zero)
                               for v in (ee, mumu, tautau))
    emu_c   = torch.complex(_rt(emu_re),   _rt(emu_im))
    etau_c  = torch.complex(_rt(etau_re),  _rt(etau_im))
    mutau_c = torch.complex(_rt(mutau_re), _rt(mutau_im))

    row0 = torch.stack((ee_c, emu_c, etau_c))
    row1 = torch.stack((emu_c.conj(), mumu_c, mutau_c))
    row2 = torch.stack((etau_c.conj(), mutau_c.conj(), tautau_c))
    return torch.stack((row0, row1, row2)).to(dtype=cdtype)
\end{lstlisting}
Every one of the eighteen scalar fields therefore accepts either a plain
Python \code{float} (producing a non-grad-tracking leaf tensor, exactly as
before) \emph{or} a differentiable \code{torch.Tensor} with
\code{requires\_grad=True}, in which case the gradient reaches
\code{epsilon}/\code{epsilon\_n} intact. \code{device=None} resolves to
CPU and every field is moved (differentiably, via \code{.to()}) onto that
single resolved device, so mixing a GPU-resident coupling with plain-float
ones never produces a mixed-device tensor. This is what lets
\pyclass{inference.solar\_model.SolarNSIOscillationModel}
(\cref{sec:impl-autograd}) fit \code{eps\_ee} by gradient descent while
building an ordinary \pyclass{NSIConfig}, with no separate raw-tensor
escape hatch needed for that case.
\end{implbox}

\begin{notebox}
\textbf{Identity-based equality.} \pyclass{NSIConfig} is declared with
\code{@dataclass(frozen=True, eq=False)}: \code{\_\_eq\_\_}/\code{\_\_hash\_\_}
are inherited from \code{object} (identity-based) rather than
auto-generated from the fields. A field-tuple-based \code{\_\_eq\_\_} would
be unsound here for two independent reasons: (i) comparing
\code{torch.Tensor} fields with \code{==} does not reliably reduce to a
plain boolean, and (ii) \pyfunc{NSIConfig.from\_raw\_epsilon}
(below) leaves the eighteen scalar fields at their \code{0.0} defaults
while setting \code{epsilon}/\code{epsilon\_n} directly, so a
field-tuple-based comparison would make a config carrying a non-zero raw
$\epsilon$ compare equal to the plain Standard-Model-limit
\pyclass{NSIConfig()} -- a real defect discovered and fixed during
external review of this extension. Use
\code{torch.equal(a.epsilon, b.epsilon)} directly for value comparison.
\end{notebox}

\begin{implbox}[title={Preset construction and the raw-epsilon escape hatch}]
\begin{itemize}
  \item \pyfunc{NSIConfig.from\_preset(name)} builds a named configuration
    from \code{tpeanuts.config.presets.NSI\_PRESETS} (\cref{ch:nsi-observables}).
  \item \pyfunc{NSIConfig.from\_raw\_epsilon(epsilon, *, epsilon\_n=None)} is
    the escape hatch for an $\epsilon$ (and, optionally, $\epsilon^n$) that
    does not decompose into the eighteen-scalar parametrization -- e.g.\ a
    precomputed matrix. It validates that each tensor supplied has at least
    two dimensions, is square, is finite, and is Hermitian within a
    dtype-scaled tolerance (raising \code{ValueError} otherwise), and, when
    both are supplied, that $\epsilon$ and $\epsilon^n$ share the same
    active-flavour block size. \textbf{A config built this way must not be
    passed through \code{dataclasses.replace()}}: since its scalar fields
    stay at their \code{0.0} defaults, \code{\_\_post\_init\_\_} would
    silently recompute \code{epsilon}/\code{epsilon\_n} back to zero on the
    next \code{replace()} call, discarding the raw tensor (and any
    gradient connected to it) with no warning. Build a fresh config with
    \code{from\_raw\_epsilon} again instead.
\end{itemize}
\end{implbox}

\begin{implbox}[title={Convenience properties}]
\begin{itemize}
  \item \pyfunc{is\_sm\_limit}: \code{True} iff \emph{both} \code{epsilon}
    and \code{epsilon\_n} are exactly the zero matrix.
  \item \pyfunc{has\_cp\_violation}: \code{True} iff any off-diagonal entry
    of \emph{either} matrix has a non-zero imaginary part.
  \item \pyfunc{has\_neutron\_coupling}: \code{True} iff \code{epsilon\_n} is
    not exactly zero. This is the single flag every downstream consumer in
    this chapter uses to decide whether the composition term
    (\cref{eq:epsilon-composition}) needs to be evaluated at all -- it costs
    nothing extra when \code{False} (the default), so every pre-existing
    scenario is bit-identical to before this extension existed.
  \item \pyfunc{eps\_emu}/\pyfunc{eps\_etau}/\pyfunc{eps\_mutau} (and their
    \code{\_n} counterparts) return the corresponding entry as a plain
    Python \code{complex}, via a helper that detaches a tensor field first
    -- these are non-differentiable display/convenience accessors, not part
    of the gradient path.
\end{itemize}
\end{implbox}

\section{Embedding $\epsilon$/$\epsilon^n$ for an arbitrary flavour count}

\begin{implbox}
\pyfunc{epsilon\_tensor(n\_flavours=...)} and its mirror
\pyfunc{epsilon\_n\_tensor(n\_flavours=...)} embed the $3\times3$ active-flavour
block in the top-left corner of an $n_{\rm flavours}\times n_{\rm flavours}$
zero matrix when $n_{\rm flavours}>3$ (the 3+1 sterile case): the sterile row
and column carry no NSI coupling of either kind, since the sterile state is
an $\mathrm{SU}(2)_L\times\mathrm{U}(1)_Y$ singlet and NSI operators are
built from Standard-Model-charged fermion currents. Both methods share a
single embedding implementation (\code{\_embed\_active}), parametrised only
by which of \code{self.epsilon}/\code{self.epsilon\_n} to embed, so the two
code paths cannot silently diverge.
\end{implbox}

\section{Hamiltonian assembly: the three-piece decomposition}
\label{sec:impl-hamiltonian}

\modulehead{core/common/hamiltonian.py}{\pyfunc{hamiltonian\_matter\_reduced}
-- the single function assembling the matter Hamiltonian for every scenario.}

\begin{theorybox}
\pyfunc{hamiltonian\_matter\_reduced} builds the flavour-basis matter term as
up to \emph{three} independent, additive pieces, then rotates the whole sum
into the reduced basis via $O^\dagger(\cdot)O$ (\S1.1, \S1.2):
\begin{equation}
  \keyresult{
    H_{\rm mat}^{\rm flavour} =
    \underbrace{V_{CC}(n_e)\bigl(\diag(1,0,\dots,0)+\epsilon\bigr)}_{\text{piece 1: NSI, }\epsilon}
    \;+\;
    \underbrace{V_{NC}(n_n)\,\diag(0,\dots,0,-1)}_{\text{piece 2: sterile NC}}
    \;+\;
    \underbrace{V_{CC}(n_n)\,\epsilon^n}_{\text{piece 3: NSI composition}}
  },
  \label{eq:three-piece}
\end{equation}
with $V_{CC}(n)\equiv\pm\sqrt2\,G_F\,n\,L_{\rm scale}$ the same charged-current
potential function of \S1.1, evaluated once on the electron density and once
(for pieces 2 and 3) on the neutron density, and $V_{NC}(n)$ the neutral-current
potential of \S1.3. Piece 3 is exactly the composition term of
\cref{eq:epsilon-composition},
\begin{equation}
  V_{CC}(n_n)\,\epsilon^n = L_{\rm scale}\sqrt2\,G_F\,n_n\,\epsilon^n
  = \Bigl(L_{\rm scale}\sqrt2\,G_F\,n_e\Bigr)\frac{n_n}{n_e}\,\epsilon^n
  = V_{CC}(n_e)\,\frac{n_n}{n_e}\,\epsilon^n,
\end{equation}
so re-evaluating the \emph{same} \pyfunc{matter\_potential\_cc} function on
the neutron density instead of the electron density is algebraically
identical to the ratio form of \cref{eq:epsilon-composition}, without ever
dividing by $n_e$ -- avoiding a division that would be singular in vacuum
($n_e=n_n=0$).
\end{theorybox}

\begin{implbox}
\begin{itemize}
  \item \code{include\_nc = n\_n\_mol\_cm3 is not None and n\_flavours == 4}:
    piece 2 is active only for the 3+1 sterile extension, and only when the
    caller supplies a neutron density; omitting it recovers the CC-only
    sterile Hamiltonian used throughout earlier phases of this project,
    exactly as before this extension existed.
  \item \code{needs\_eps\_n = oscillation.BSM\_extension\_NSI and
    oscillation.nsi.has\_neutron\_coupling}: piece 3 is active for
    \emph{any} flavour count, driven purely by whether $\epsilon^n\neq0$ --
    not by whether the sterile extension is also active.
  \item \textbf{Piece 3 is not optional once requested.} If
    \code{needs\_eps\_n} is \code{True} but \code{n\_n\_mol\_cm3} was not
    supplied, the function raises \code{ValueError} rather than silently
    falling back to the $\epsilon$-only Hamiltonian. This is a deliberate
    asymmetry with piece 2: omitting \code{n\_n\_mol\_cm3} for the sterile
    NC term is itself a well-defined, physically exact choice (``CC-only
    sterile propagation''), whereas a non-zero $\epsilon^n$ is an explicit
    user request that the function must not silently downgrade.
  \item The direction matrices $P_{cc}$, $P_{nc}$, $P_{\epsilon^n}$ (the
    $O^\dagger(\cdot)O$-rotated, $V$-independent counterparts of the three
    pieces above) are built once by the shared helper
    \pyfunc{\_matter\_direction\_reduced}, reused identically by both the
    exact Hamiltonian builder here and the perturbative segment evolutor
    (\cref{sec:impl-perturbative}).
\end{itemize}
\end{implbox}

\begin{notebox}
Because \pyfunc{hamiltonian\_matter\_reduced} is the single function every
other builder in this chapter delegates to, this \code{ValueError} is a
defense-in-depth guarantee: any current or future caller anywhere in the
pipeline that forgets to supply the neutron density for a non-zero
$\epsilon^n$ configuration fails loudly at the point of Hamiltonian
assembly, rather than silently returning a physically incomplete result.
\end{notebox}

\section{The perturbative segment evolutor}
\label{sec:impl-perturbative}

\modulehead{core/perturbative/evolutor.py}{\pyfunc{evolutor\_perturbative\_segment}
-- first-order perturbative evolution for one piecewise-polynomial density
segment (used by \texttt{medium.earth} and \texttt{medium.atmosphere}).}

\begin{theorybox}
For a segment over which the electron/neutron density varies, the
zeroth-order evolutor $U_0$ uses the segment-\emph{average} potential;
the first-order correction $U_1$ is built from the spectral projectors
$M_a$ of the average Hamiltonian and the oscillatory integral of the
\emph{residual} density (density minus its segment average) against the
Hamiltonian eigenvalue differences,
\begin{equation}
  U_1 = -i\sum_{a,b} I_{ab}\,\bigl(M_a\,P\,M_b\bigr),
  \qquad
  I_{ab} = \int_{x_1}^{x_2}\delta n(x)\,e^{-i\lambda_a x+i\lambda_b x}\,dx,
\end{equation}
where $P$ is whichever position-independent direction matrix the term in
question needs (\cref{eq:three-piece}).
\end{theorybox}

\begin{implbox}
\pyfunc{evolutor\_first\_order} accepts three independent, optional direction
matrices, mirroring \cref{eq:three-piece} exactly: \code{P} (piece 1),
\code{P\_nc} (piece 2), and \code{P\_eps\_n} (piece 3). The key implementation
observation is that pieces 2 and 3 need the \emph{same} underlying
oscillatory integral -- both are built from the neutron-density residual,
\code{profile\_model.residual\_integral\_neutron(la, lb)} -- and differ only
in which potential-conversion function is applied to it afterwards
(\pyfunc{matter\_potential\_nc} for piece 2, \pyfunc{matter\_potential\_cc}
for piece 3, exactly matching \cref{eq:three-piece}'s $V_{NC}(n_n)$ versus
$V_{CC}(n_n)$) and which direction matrix they sandwich. \code{P\_eps\_n}
therefore required \emph{no new residual-integral machinery}: it reuses
\code{residual\_integral\_neutron} exactly as \code{P\_nc} already did.
\end{implbox}

\begin{implbox}[title={Zeroth-order composition term and the neutron-coupling guard}]
\pyfunc{evolutor\_perturbative\_segment} computes
\code{include\_eps\_n = oscillation.BSM\_extension\_NSI and
oscillation.nsi.has\_neutron\_coupling} and, when true, builds the
zeroth-order piece-3 contribution from \code{profile\_model.average\_n}
(the segment-average neutron density) via \pyfunc{matter\_potential\_cc} --
\emph{not} from \code{profile\_model.potential\_n}, which is already
converted with the \emph{neutral-current} prefactor for piece 2's own use
and would apply the wrong prefactor if reused directly for piece 3. As in
\cref{sec:impl-hamiltonian}, a non-zero $\epsilon^n$ with a
\code{profile\_model} that was not built with neutron-density coefficients
(\code{average\_n}/\code{potential\_n}/\code{residual\_integral\_neutron}
unavailable) raises \code{ValueError} rather than silently ignoring the
request -- mirroring, and reusing the same underlying neutron-density data
as, the pre-existing \code{include\_matter\_nc} guard for piece 2.
\end{implbox}

\section{Pipeline wiring: fetching neutron density for the right reason}
\label{sec:impl-pipeline}

\begin{notebox}
\textbf{The subtlety this section documents.} Piece 2 (sterile NC) and
piece 3 (NSI composition) both need local neutron-density data, but they
are triggered by \emph{independent} conditions: piece 2 by the
sterile-specific, tri-state \code{include\_matter\_nc} policy
(\pyfunc{core.common.oscillation.resolve\_include\_matter\_nc}, unchanged
by this extension), piece 3 automatically by
\code{oscillation.nsi.has\_neutron\_coupling}, with no separate flag --
exactly like $\epsilon$ itself needs no ``enable NSI'' flag beyond
\code{oscillation.nsi} being set. Every \texttt{medium.*} pipeline that
decides \emph{whether to sample/fetch} raw neutron-density data at all
(as opposed to whether to \emph{apply} the sterile NC term, which is
already independently gated downstream) must therefore OR the two
conditions together, via the shared helper
\begin{lstlisting}[style=tpython]
def oscillation_needs_neutron_composition(oscillation):
    return oscillation.BSM_extension_NSI and oscillation.nsi.has_neutron_coupling
\end{lstlisting}
in \texttt{core.common.oscillation}.
\end{notebox}

\begin{implbox}
\begin{itemize}
  \item \textbf{Earth.} \pyclass{EarthProfile} bakes its neutron-density
    coefficients in at \emph{construction} time
    (\code{profile\_perturbative\_kwargs={'include\_neutron': True}}),
    independently of any per-call flag -- so the exact
    (\pyfunc{evolutor\_numerical}-based) and perturbative
    (\pyfunc{evolutor\_perturbative\_segment}-based) Earth paths both
    automatically get composition-term support for free once the profile was
    built with neutron data, requiring \emph{no} pipeline-level change. The
    one exception is \pyfunc{earth\_probability\_state\_numerical}, which
    (unlike the perturbative Earth path) samples \code{n\_n} itself from the
    profile and previously did so only when \code{include\_matter\_nc} was
    true; it now also samples it when
    \pyfunc{oscillation\_needs\_neutron\_composition} is true, raising a
    dedicated \code{ValueError} if the profile lacks neutron data.
  \item \textbf{Atmosphere.} Both \pyfunc{atmosphere\_evolutor\_numerical}
    and \pyfunc{atmosphere\_evolutor\_analytical} decide whether to
    sample/fit neutron density \emph{per call} (unlike Earth); both now OR
    in \pyfunc{oscillation\_needs\_neutron\_composition} alongside the
    resolved \code{include\_matter\_nc}.
  \item \textbf{Sun.} \pyfunc{solar\_evolutor\_numerical\_history} (feeding
    \pyfunc{core.numerical.evolutor\_numerical}) and
    \pyfunc{solar\_probability\_mass}'s \code{method="adiabatic\_exact"}
    branch (feeding \pyfunc{mass\_weights\_adiabatic\_exact}, which itself
    calls \pyfunc{hamiltonian\_flavour}) both apply the same OR condition
    before sampling \code{medium.density\_n}. A related, more subtle fix was
    needed inside \pyfunc{mass\_weights\_adiabatic\_exact} itself: its
    fixed-permutation bookkeeping (relabelling \code{torch.linalg.eigh}'s
    ascending-eigenvalue column order by physical vacuum mass index, needed
    because that order disagrees with the physical one under inverted
    ordering) diagonalises an auxiliary \emph{vacuum-limit} Hamiltonian at
    $n_e=n_n=0$; that internal call used to pass \code{n\_n\_mol\_cm3=None}
    (physically inert either way, since $V_{CC}(0)=0$), which now trips the
    \S\ref{sec:impl-hamiltonian} guard whenever $\epsilon^n\neq0$. It now
    passes an explicit zero tensor instead of \code{None} -- numerically
    identical, but compatible with the new guard.
\end{itemize}
\end{implbox}

\begin{notebox}
\textbf{method="adiabatic\_approximated" remains Standard-Model only.} The
closed-form two-level matter-mixing-angle formulas underlying that method
have no NSI generalisation at all (not even for $\epsilon^n=0$); NSI of
either kind is rejected there with \code{ValueError} regardless of
composition, forcing \code{method="adiabatic\_exact"} or
\code{method="numerical"} for any NSI scenario in the Sun.
\end{notebox}

\section{Not yet supported}

\begin{notebox}
The composition term (piece 3) has \emph{no} generalisation implemented for
\pyfunc{medium.solar.adiabatic.mass\_weights\_adiabatic\_approximated} or
its Landau-Zener correction, for the same structural reason $\epsilon$
itself has none there (\S\ref{sec:impl-pipeline}). It is fully implemented
for the exact evolutor path (any medium) and for the perturbative segment
evolutor (Earth, atmosphere); the only propagation \emph{method} left
without composition-term support is the closed-form adiabatic
approximation, which was already NSI-incompatible before this extension.
\end{notebox}

\section{Autograd: fitting $\epsilon$ by gradient descent}
\label{sec:impl-autograd}

\modulehead{inference/solar\_model.py}{\pyclass{SolarNSIOscillationModel} --
a differentiable solar $P_{ee}(E)$ model with a free diagonal NSI coupling.}

\begin{implbox}
\pyclass{SolarNSIOscillationModel.oscillation(theta)} builds a fresh
\pyclass{NSIConfig} directly from the free-parameter tensor,
\begin{lstlisting}[style=tpython]
nsi = NSIConfig(
    eps_ee=values["eps_ee"],          # a torch.Tensor, requires_grad=True
    device=self.context.device,
    real_dtype=self.context.dtype,
)
\end{lstlisting}
relying on \cref{sec:impl-nsiconfig}'s differentiable
\pyfunc{\_hermitian\_3x3} to keep \code{eps\_ee}'s gradient connected all the
way through \code{epsilon}, \pyfunc{hamiltonian\_flavour}, and
\pyfunc{mass\_weights\_adiabatic\_exact}'s
\code{torch.linalg.eigh} diagonalisation, to the predicted $P_{ee}$. This
constructor call is what \pyfunc{NSIConfig.from\_raw\_epsilon} (a one-hot
$\times$ \code{eps\_ee} matrix multiplication, in an earlier version of this
model) existed to work around before \cref{sec:impl-nsiconfig}'s rewrite --
the differentiable scalar-field path removes the need for that workaround
for any coupling that decomposes into the ordinary
\code{eps\_*}/\code{eps\_*\_n} parametrization, which a single diagonal
coupling like \code{eps\_ee} always does.
\pyfunc{inference.fit.fit\_lbfgs} then differentiates through the whole
chain with \code{torch.optim.LBFGS}, exactly as it already does for the
purely-Standard-Model free parameters (\code{theta12}, \code{DeltamSq21},
\dots); see \cref{sec:results-borexino-nsi} for the resulting fit.
\end{implbox}

\section{Summary: modules touched by this extension}

\begin{center}
\begin{tabular}{@{}ll@{}}
\toprule
Module & Role \\
\midrule
\pkgpath{core/BSM/bsm\_nsi.py} & \pyclass{NSIConfig}: parameters, differentiable construction, embedding \\
\pkgpath{core/common/hamiltonian.py} & \pyfunc{hamiltonian\_matter\_reduced}: the three-piece decomposition \\
\pkgpath{core/common/oscillation.py} & \pyfunc{oscillation\_needs\_neutron\_composition} \\
\pkgpath{core/numerical/evolutor.py} & Exact matrix-exponential evolution; consumes \code{n\_n\_mol\_cm3} generically \\
\pkgpath{core/perturbative/evolutor.py} & \pyfunc{evolutor\_perturbative\_segment}: \code{P\_eps\_n} sandwich term \\
\pkgpath{medium/earth/probability.py} & Neutron-density sampling for \pyfunc{earth\_probability\_state\_numerical} \\
\pkgpath{medium/atmosphere/evolutor.py} & Neutron-density sampling for both atmosphere propagation methods \\
\pkgpath{medium/solar/evolutor.py} & Neutron-density sampling for the numerical solar evolutor \\
\pkgpath{medium/solar/probability.py} & Neutron-density sampling for \code{method="adiabatic\_exact"} \\
\pkgpath{medium/solar/adiabatic.py} & Vacuum-limit permutation fix in \pyfunc{mass\_weights\_adiabatic\_exact} \\
\pkgpath{inference/solar\_model.py} & Differentiable \code{eps\_ee} fit via \pyclass{SolarNSIOscillationModel} \\
\bottomrule
\end{tabular}
\end{center}
